\documentclass{article}
\usepackage{comment}
\usepackage[hypertexnames=false,colorlinks=true,breaklinks]{hyperref}
\usepackage{url}
\pagestyle{empty}		% no page numbers, thanks.
\addtolength{\topmargin}{-0.5in}% repairing LaTeX's huge margins...
% \setlength{\oddsidemargin}{0.5in}
\setlength{\oddsidemargin}{0.5in}
\addtolength{\textwidth}{1.5in}	% and here...
\addtolength{\marginparsep}{-0.5in}
\addtolength{\marginparwidth}{0.5in}
\setlength{\parindent}{0mm}	% indented paragraphs just don't make it
			   	% on a resume.  (My opinion; if you
			   	% disagree, give this whatever value you
			   	% want.)
\setlength{\textheight}{9.0in}	% more margin hacking.  I have a large
%\setlength{\textheight}{10.2in}	% more margin hacking.  I have a large
			    	% resume to fit on one page.

\begin{document}
\reversemarginpar		% this puts the margin notes on the
				% left-hand side of the resume.

{\huge{Joaqu\'{i}n Rapela}}	% make the name stand out

				% address stuff.  The \hspace*{\fill} is
				% necessary to get things laid out
				% correctly.  The sizes fed to \vspace are
				% variable, depending on how much
				% whitespace you want, and how much
				% stuff you're trying to fit on a page....
% \vspace{.1in}
\renewcommand{\labelitemi}{$\star$}
\textit{\textbf{Address}}: Gatsby Computational Neuroscience Unit. University
College London. 25 Howland St, Fitzrovia, London W1T 4JG, United Kingdom.
\textit{\textbf{Phone}}:~44 7432049398
\textit{\textbf{email}}:~j.rapela@ucl.ac.uk
\textit{\textbf{www}}:~http://www.gatsby.ucl.ac.uk/$\sim$rapela,
\textit{\textbf{github}}:~https://github.com/joacorapela,
\textit{\textbf{last update}}:~\today

%\par
% \vspace{2\baselineskip}
%\vspace{\baselineskip}
%A full-time position in software engineering or  \marginpar{{\bf Objective}}
%user interface design.
				% if the \marginpar is at the beginning
				% of the line, it loses.  *sigh*.  This
				% is a typical entry in a resume.   (I'm
				% defining an entry as something
				% important enough to have a marginal note
				% setting it off.)  Give it some space
				% with \vspace, use \par to break
				% between paragraphs, enter your text,
				% and place the \marginpar at the end.
\par
\vspace{2\baselineskip}
{\bf University of Southern California} \marginpar{{\bf Education}}

\par
\vspace{.05in}
{\bf Department of Electrical Engineering} \hspace*{\fill} 2003---2010
\par

\par
PhD in Electrical Engineering.
\par
Advisors: Dr. Norberto M. Grzywacz (Neuroscience) and Dr. Jerry M. Mendel (Signal Processing)
\par
\vspace{.05in}
{\bf Neuroscience Graduate Program} \hspace*{\fill} 2001---2003
\par

\par
Studies in Neuroscience.

\vspace{.05in}
{\bf Department of Electrical Engineering} \hspace*{\fill} 2000---2003

\par
MS in Electrical Engineering.

\par
\vspace{.1in}
{\bf Universidad de Buenos Aires}
\par
``Licenciatura (equivalent to MS)'' in Computer Science.
\hfill 1995---1998
\par
BS in Computer Science.\hfill 1992---1995

\par
\vspace{2\baselineskip}
\par
Aishah Qureshi, Campagner Dario, Mitra Javadzadeh, \textbf{Rapela J.} (2022)
\marginpar{{\bf Publications}}
\textit{Tracking Freely Moving Mice Using Computer Vision, Statistical Inference
and Statistical Learning Techniques}.
44th Annual International Conference of the IEEE Engineering in Medicine \&
Biology Society (EMBC)
\href{http://www.gatsby.ucl.ac.uk/~rapela/posters/EMBC22_1882_MS.pdf}{one-page
paper},
\href{http://www.gatsby.ucl.ac.uk/~rapela/posters/EMBC22_1882_poster.pdf}{presented
poster}.


\par
Javadzadeh M, \textbf{Rapela J.}, Sahani M., Hofer S.B. (2022, selected for oral presentation)
\textit{Dynamic causal communication channels between neocortical areas}.
CoSyNe.
\href{http://www.gatsby.ucl.ac.uk/~rapela/mitra/abstracts/cosyne21/cosyne_abstract_Javadzadeh.pdf}{abstract}.

\par
\textbf{Rapela J.}, (2022)
\textit{Matrix differentials simplify the calculation of derivatives with
respect to matrices.}.
\href{http://www.gatsby.ucl.ac.uk/~rapela/papers/ppcaLLmatrixDerivatives.pdf}{technical report}.

\par
\textbf{Rapela J.}, Todorov D. (accepted)
\textit{Hidden Markov models applied to LFPs from layer two and three of human
cortex reveal highly stereotypical discrete states in epileptic seizures
separated by more than an hour}. 2019 41st Annual International Conference of
the IEEE Engineering in Medicine and Biology Society (EMBC).
\href{http://www.gatsby.ucl.ac.uk/~rapela/papers/Rapela_and_Todorov_EMBC_2019.pdf}{article},
\href{http://www.gatsby.ucl.ac.uk/~rapela/papers/embc19_hmm_supplemental.pdf}{supplemental}.

\par
\textbf{Rapela J.}, Proix T., Todorov D., Truccolo W. (2019, selected for oral presentation)
\textit{Uncovering low-dimensional structure in high-dimensional
representations of long-term recordings in people with epilepsy}. 2019 41st Annual International Conference of
the IEEE Engineering in Medicine and Biology Society (EMBC).
\href{http://www.gatsby.ucl.ac.uk/~rapela/papers/EMBC19_Rapela_et_al_IEEE_EMBC_2019.pdf}{IEEE EBMC},
\href{https://embs.papercept.net/conferences/conferences/EMBC19/program/EMBC19_ContentListWeb_3.html#thc02_02}{preprint}.

\par
\textbf{Rapela J.}, Westerfield M., Townsend J. (2018)
\textit{A new foreperiod effect on single-trial phase coherence. Part I:
existence and relevance}. Neural Computation 30(9):2348-2383.
\href{https://www.mitpressjournals.org/doi/abs/10.1162/neco_a_01109}{Neural Computation},
\href{http://www.gatsby.ucl.ac.uk/~rapela/papers/NECO-09-17-2964-main.pdf}{preprint},
\href{http://www.gatsby.ucl.ac.uk/~rapela/papers/NECO-09-17-2964-supplemental.pdf}{preprint supplemental}
(3 citations).

\par
\textbf{Rapela J.} (2018)
\textit{Traveling waves appear and disappear in unison with produced speech}.
\href{https://arxiv.org/abs/1806.09559}{preprint} (4 citations).

\par
\textbf{Rapela J.} (2017)
\textit{Rhythmic production of consonant-vowel syllables synchronizes traveling waves in speech-processing brain regions}.
\href{https://arxiv.org/abs/1705.01615}{preprint} (2 citations).

\par
\textbf{Rapela J.} (2016)
\textit{Entrainment of traveling waves to rhythmic motor acts}.\linebreak
\href{http://arxiv.org/abs/1606.02372}{preprint} (2 citations).

\par
\textbf{Rapela J.}, Kostuk M, Rowat, P.F., Mullen, T., Chang E.F., Bouchard K.
(2015)
\textit{Modeling neural activity at the ensemble level}.
\href{http://arxiv.org/abs/1505.00041}{preprint}.

\par
\textbf{Rapela J.}, Gramann K., Westerfield M., Townsend J., \& Makeig S.
(2012).  \textit{Brain oscillations in Switching vs. Focusing audio-visual
attention}.  Proceedings of the 34th Annual International Conference of the
IEEE EMBS, San Diego, California. (15 citations)

\par
\textbf{Rapela J.}, Tsong-Yan L., Westerfield M., Jung T.P. \& Townsend J.,
(2012).  \textit{Assisting autistic children with wireless EOG technology}.
Proceedings of the 34th Annual International Conference of the IEEE EMBS, San
Diego, California. (14 citations)

\par
\textbf{Rapela J.}, Felsen G., Touryan J., Mendel J.M., \& Grzywacz N.M. (2010).
\textit{ePPR: a new strategy for the characterization of sensory cells from input/output data}.
Network: Computation in Neural Systems. 21(1-2): 35-90. (15 citations)

\par
\textbf{Rapela J.}, Mendel J.M., \& Grzywacz N.M. (2006).
\textit{Estimating nonlinear receptive fields from natural images}.
Journal of Vision 6(4), 441-474. (31 citations)

\par
Shattuck D., \textbf{Rapela, J.}, Asma E., Chatziioannou A., Qi J., \& Leahy R. (2002).
\textit{Internet2-based 3D PET Image Reconstruction using a PC Cluster}.
Physics in Medicine and Biology 47, 2785-2795. (57 citations)

\par
\textbf{Rapela J.} (2001).
\textit{Automatically Combining Ranking Heuristics for HTML Documents}.
Proceedings of the Third International Workshop on Web Information and Data
Management, 61-67, New York, NY: ACM Press. (18 citations)

\par
\vspace{2\baselineskip}
{\bf Gatsby Computational Neuroscience Unit}
\marginpar{{\bf Academic}}
\par
{\bf University College London} \hspace*{\fill} June 2019---
\marginpar{{\bf Experience}}
\par
{\bf Position}: Research Engineer Fellow.
\par

\begin{list}{\labelitemi}{\leftmargin=1em}

\item Distributes, performs quality control and enhances advanced statistical
    neuroscience algorithms developed at the Gatsby Computational Neuroscience
    Unit (e.g., \href{http://www.github.com/joacorapela/svGPFA}{Sparse
    Variational Gaussian Processes Factor Analysis}) (keywords: sparse
    variational learning, Gaussian processes, nonlinear optimization,
    unsupervised learning, Python, continuous integrations, Pytorch, SciPy,
    Plotly).

\item Collaborates with experimental neuroscientists at the Sainsbury Wellcome
    Centre on the use of advanced statistical algorithms for the
    characterization of neural populations electrophysiological recordings
    (keywords:
    \href{https://github.com/joacorapela/ldsForNeuralPopulation}{linear
    dynamical systems with external inputs}, expectation maximization).

\item Contributes to the development of software infrastructure to record,
    store, manage and display large-scale continuous (24/7) behavioral and
    electrophysiological recordings from foraging mice (keywords: foraging,
    dahsboard, interactive visualization, data aquistion, experimental
    control, Python, software development best practices).

\item As part of the Simons Collaborative on the Global Brain (SCGB)
    undergraduate research fellowship (SURF), with Dr.~Mitra Javadzadeh No
    (Sainsbury Wellcome Centre), co-mentors Ms.~Aishah Qureshi, undergraduate
    student from Queen Mary University of London
    (\href{https://www.simonsfoundation.org/people/joaquin-rapela/}{more
    info}).

\item Trains members of the Gatsby Computational Neuroscience Unit and the
    Sainsbury Wellcome Centre in open-science tools (e.g.,
    \href{https://sainsburywellcomecentre.github.io/pystarters/events/intro_nov_2019.html}{Introduction
    to Python and Open-Science Tools},
    \href{https://github.com/joacorapela/pyClubCI}{Hands-On Tutorial on Code
    Testing for Researchers}).

\item Works with members of the Sainsbury Wellcome Centre and the Gatsby
    Computational Neuroscience Unit to improve the research culture of our
    workplace (\href{https://sainsburywellcomecentre.github.io/RCWG/}{more info}).

\end{list}

\par
\vspace{2\baselineskip}
{\bf The Truccolo Lab for Computational Neuroscience}
\par
{\bf Brown University} \hspace*{\fill} August 2017---January 2019
\par
{\bf Position}: Research Associate.
\par

\begin{list}{\labelitemi}{\leftmargin=1em}

\item Discovered stereotypical discrete states in pre-ictal, ictal and
post-ictal periods of seizures separated by more than one hour. This finding
was obtained using micro-electrode array recordings from cortical layers two
and three of persons with epilepsy (Rapela and Todorov, 2019) (keyword:
epilepsy, unsupervised time-series modeling, hidden Markov models).

\item Found two-dimensional descriptors of high-dimensional representations of
electrophysiological neural recordings from persons with epilepsy that
separate well inter-ictal, pre-ictal, ictal and post-ictal periods.
%
This finding was obtained using micro-electrode array recordings from cortical
layers two and three of persons with epilepsy (Rapela et al, 2019)
(keywords: epilepsy, low-dimensional manifolds, t-SNE, XGBoost, multivariate
logistic regression).

\end{list}

\par
\vspace{2\baselineskip}
{\bf Instituto en Luz Ambiente y Vision}
\par
{\bf Universidad Nacional de Tucum\'{a}n, Argentina} \hspace*{\fill} November 2013---Present
\par
{\bf Positions}: External researcher.
\par

\par
\vspace{2\baselineskip}
{\bf Swartz Center for Computational Neuroscience}
\par
{\bf University of California San Diego} \hspace*{\fill} November 2010---July 2017
\par
{\bf Positions}: Postdoctoral and visiting scholar,
\par

\begin{list}{\labelitemi}{\leftmargin=1em}

\item Characterized spatio-temporal brain dynamics related to speech
production from electrocorticographic neural recordings (Rapela 2016, 2017)
(keywords: spatio-temporal dynamics, synchronization of oscillators, phase-amplitude coupling, phase coherence)

\item Principal investigator in the project \emph{Taking the next step toward
understanding computations by neural ensembles with high resolution
neural recordings, generic data assimilation methods, and increased
computational power} funded by the Center for Brain Activity Mapping, part of
president's Obama B.R.A.I.N. initiative (Rapela et al., 2015) (keywords: population density models, dynamical systems)

\item Discovered a new effect of temporal expectation on the single-trial
phase coherence of the electroencephalogram (EEG) using a variational Bayes
linear regression method. (Rapela et al., 2018) (keywords: attention,
expectation, EEG, time-frequency analysis, single trial
analysis,variational-bayes linear regression, pattern recognition and machine
learning)

\item Found that switching attention between the visual and auditory modality
generate transient arousal of attention in both modalities (Rapela, Gramman et
al., 2012) (keywords: attention switch, EEGLAB, time-frequency analysis)

\item Developed a computer game controlled online by the eye movements of the
player, recorded using electrooculography, and quantified the speed and
accuracy of attention-orienting eye movements (Rapela, Line et al., 2012)
(keywords: eye movements, EEG, EOG)

\end{list}

\par
\vspace{2\baselineskip}
{\bf University of Southern California}
\par
{\bf Position}: Research Assistant. \hspace*{\fill} 2001---2010
\par
{\bf Advisors}: Dr. Norberto M. Grzywacz. Director, Neuroscience Graduate
Program. Professor, Department of Biomedical Engineering. Dr. Jerry Mendel.
Professor, Department of Electrical Engineering.

\begin{list}{\labelitemi}{\leftmargin=1em}
\item Investigated Bayesian models to evaluate the optimality of retinal computations (keywords: Bayesian statistics, optimal neural codes).
\item Developed the Volterra Relevant Space Technique (VRST) for the
estimation of spatial Volterra models of visual cells (Rapela et al., 2006) (keywords: Volterra models, dimensionality reduction, nonlinear systems, regression analysis)
\item Developed the extended Projection Pursuit Regression (ePPR) algorithm
for the spatio-temporal characterization of visual cells from input/output
data (Rapela et al., 2010) (keywords: projection pursuit regression, nonlinear
optimization, dimensionality reduction)

\end{list}

\par
\vspace{.1in}
{\bf University of Southern California}
\par
{\bf Position}: Research Assistant. \hspace*{\fill} August 2000---August 2001
\par
{\bf Advisor}: Dr. Richard M. Leahy. Professor, Department of Electrical
Engineering.

\begin{list}{\labelitemi}{\leftmargin=1em}

\item Participated in the development of a distributed computing application
for MAP reconstructions of 3D PET data, which was accessed remotely using a
Java-based interface (Shattuck, Rapela, et al., 2002)

\item Research on methods for MEG/EEG source localization.
\end{list}

\par
\vspace{2\baselineskip}
{\bf University College London} \hspace*{\fill} November 2019, May 2021
\marginpar{{\bf Teaching}}

\par
{\bf Course}: PyStarters
\marginpar{{\bf Experience}}

\par
{\bf Role}: Instructor on Python programming.

\par
\vspace{.1in}
{\bf Brown University} \hspace*{\fill} September 2017---December 2017
\par
{\bf Course}: NEUR2110 Statistical Neuroscience
\par
{\bf Role}: Teaching assistant

\par
\vspace{.1in}
{\bf University of Southern California} \hspace*{\fill} January 2009---May 2009
\par
{\bf Course}: EE 364: Introduction to Probability and Statistics for Electrical Engineering and Computer Science. Evaluations available upon request.
\par
{\bf Role}: Discussion section leader.

\par
\vspace{.1in}
{\bf Universidad de Buenos Aires} \hspace*{\fill} August 1999---December 1999
\par
{\bf Course}: Numerical Linear Algebra.
\par
{\bf Role}: Teaching Assistant.

\par
\vspace{.1in}
{\bf Universidad de Buenos Aires} \hspace*{\fill} August 1998---August 1999
\par
{\bf Course}: Artificial Intelligence.
\par
{\bf Role}: Teaching Assistant.

\par
\vspace{2\baselineskip}
{\bf Vision Research, Frontiers in Neuroscience, Journal of Perceptual Imaging}
\marginpar{{\bf Service}}
\par
Ad hoc reviewer.
\par

\par
\vspace{2\baselineskip}
{\bf Center for Brain Activity Mapping}  \hspace*{\fill} 2013---2014
\marginpar{{\bf Funding }}
\par
Principal investigator, grant \texttt{2013-023CBAM}
\par
\emph{Taking the next step toward understanding computations by neural
ensembles with high resolution neural recordings, generic data assimilation
methods, and increased computational power}.

\par
\vspace{2\baselineskip}
{\bf University of Southern California}
\marginpar{{\bf Selected}}
\par
\marginpar{{\bf Courses}}
{\bf Mathematics:} Topology, Fundamental Concepts of Analysis, Real Analysis (audited), Functional Analysis (audited). {\bf Engineering:} Transform Theory for Engineers, Random Processes in Engineering, Introduction to Digital Signal Processing, Computational Solution to Optimization Problems, Information Theory, Statistics for Engineers. {\bf Neuroscience:} Control and Communication in the Nervous System, Neurobiology, Advanced Neurosciences I, Advanced Neurosciences II.\par
\par

\vspace{2\baselineskip}
{\bf IBM Almaden Research Center}
\marginpar{{\bf Industry}}
\par
{\bf Position}:  Staff Software Engineer
\hspace*{\fill} January 2000---August 2000
\marginpar{{\bf Experience}}
\par
Developed a Java interface for the Andrew File System, allowing users to
securely access the file system over the Internet using servlets. Experience
in distributed file systems.

\par
\vspace{.1in}
{\bf McLees, Argentina}
\par
{\bf Position}:  Software Engineer  \hfill March 1999---September 1999
\par
Participated in the development of a distributed object-oriented application in Java using CORBA and RMI.

\par
\vspace{.1in}
{\bf Alfanuclear, Argentina}
\par
{\bf Position}:  Software Engineer  \hfill December 1998---March 1999
\par
Medical Image software development in C++.

\par
\vspace{.1in}
{\bf Oracle, Argentina}
\par
{\bf Position}:  Software Engineer  \hfill June 1998---September 1998
\par
Constructed a web based application for purchasing gas station products through
the Internet using Oracle tools: Oracle Web Application Server, Oracle
Forms, and Oracle Designer 2000.

\par
\vspace{.1in}
{\bf IBM, Argentina}
\par
{\bf Position}:  Fellowship holder  \hfill September 1997---June 1998
\par
Technical consultant and programmer of Internet based applications.

\vspace{.1in}
{\bf IBM Almaden Research Center}
\nopagebreak
\par
{\bf Position}:  Fellowship holder  \hfill October 1996---September 1997
\nopagebreak
\par
Migrated to Java IBM's storage solution ADSM. Experience in JDBC, native
methods, user interface development in Java, and push technology.

\par
\vspace{.1in}
{\bf IBM, Argentina}
\par
{\bf Position}:  Fellowship holder  \hfill April 1994---October 1996
\par
Research on object persistence for the Object Oriented Software Group of IBM Argentina and development of object-oriented applications with VisualAge for Smalltalk.
Construction of an image processing application performing optical
character recognition on credit card receipts. Due to performance constraints
this application used concurrent programming with shared memory for
inter-process communication under AIX (IBM's UNIX).

\vspace{2\baselineskip}
R, Python, Matlab, Java, C/C++, Smalltalk, Pascal, Assembler.
\marginpar{{\bf Programming}}
\par
\marginpar{{\bf Languages}}
\par

\end{document}

